\setlength{\tabcolsep}{3pt}
\begin{longtable}{|L{1.2cm}|L{1.1cm}|C{2cm}|L{6cm}|}
\hline
{\bf{HEX}} & {\bf{DEC}} & {\bf{Signaal}} & {\bf{Omschrijving}} \\
\hline
\endfirsthead
\multicolumn{3}{l@{}}{\ldots vervolgd}\\
\hline
{\bf{HEX}} & {\bf{DEC}} & {\bf{Signaal}} & {\bf{Omschrijving}} \\
\hline
\endhead
\multicolumn{3}{l@{}}{vervolgd \ldots}\\
\endfoot
\hline
\endlastfoot
\footnotesize  00\index{\$00 (STOPTX)}\index{STOPTX} & \footnotesize 0 & STOPTX\index{STOPTX} & Stop onmiddellijk met het verzenden van het huidige Ethernet-frame.  De zender stuurt een gedeeltelijk verzonden frame, waardoor dit frame zeer waarschijnlijk verloren gaat.   \\
\hline
\footnotesize  01\index{\$01 (STARTTX)}\index{STARTTX} & \footnotesize 1 & STARTTX\index{STARTTX} & Zend pakket \\
\hline
\footnotesize  D0\index{\$D0 (RXNORMAL)}\index{RXNORMAL} & \footnotesize 208 & RXNORMAL\index{RXNORMAL} & Schakel de effecten van RXONLYONE uit \\
\hline
\footnotesize  D4\index{\$D4 (DEBUGVIC)}\index{DEBUGVIC} & \footnotesize 212 & DEBUGVIC\index{DEBUGVIC} & Selecteer de VIC-IV-foutopsporingsstroom via Ethernet wanneer \$D6E1.3 is gezet \\
\hline
\footnotesize  DC\index{\$DC (DEBUGCPU)}\index{DEBUGCPU} & \footnotesize 220 & DEBUGCPU\index{DEBUGCPU} & Selecteer de CPU-foutopsporingsstroom via Ethernet wanneer \$D6E1.3 is gezet \\
\hline
\footnotesize  DE\index{\$DE (RXONLYONE)}\index{RXONLYONE} & \footnotesize 222 & RXONLYONE\index{RXONLYONE} & Ontvang precies één Ethernet-frame en behoud alle signaaltoestanden (voor het opsporen van fouten in het Ethernet-subsysteem) \\
\hline
\footnotesize  F1\index{\$F1 (FRAME1K)}\index{FRAME1K} & \footnotesize 241 & FRAME1K\index{FRAME1K} & Selecteer frames van ~1 KiB voor de video-/CPU-foutopsporingsstroom (voor ontvangers die MTU's groter dan 2 KiB niet ondersteunen) \\
\hline
\footnotesize  F2\index{\$F2 (FRAME2K)}\index{FRAME2K} & \footnotesize 242 & FRAME2K\index{FRAME2K} & Selecteer frames van ~2 KiB voor de video-/CPU-foutopsporingsstroom, voor optimale prestaties. \\
\hline
\end{longtable}
